\documentclass[11pt,a4paper]{article}
\usepackage[utf8]{inputenc}
\usepackage[T1]{fontenc}
\usepackage[ngerman]{babel}
\usepackage[margin=2.4cm]{geometry}
\usepackage{booktabs}
\usepackage{longtable}
\usepackage{xcolor}
\usepackage[colorlinks=true,urlcolor=blue!60!black,linkcolor=black]{hyperref}
\usepackage{parskip}
\usepackage{enumitem}
\setlist{nosep,leftmargin=*}

\newcommand{\code}[1]{\texttt{\small #1}}
\newcommand{\ok}{\textcolor{green!45!black}{\textbf{erledigt}}}
\newcommand{\offen}{\textcolor{orange!75!black}{\textbf{offen}}}
\newcommand{\upstream}{\textcolor{blue!60!black}{\textbf{upstream}}}

\title{AWI-ESM3-4-2-veg-HR, CMIP7 piControl\\[2mm]
\large Antwort auf das DKRZ-Review von cli108\\ und Ergebnisse der anschlie\ss enden eigenen Nachsuche}
\author{Jan Streffing, AWI}
\date{8. August 2026}

\begin{document}
% babel-ngerman macht " aktiv, was in Code-Beispielen zuschlaegt:
% operation="instant" wurde als operation=ïnstant" gesetzt. Die Shorthands
% werden erst bei \begin{document} scharfgeschaltet, also muss das hierhin.
\shorthandoff{"}
\maketitle

\section{Kurzfassung}

Sie haben uns zu \code{cli108} (Version \code{v20260804}) zwei Runden geschickt.
Beide sind jetzt abgearbeitet. Der Verifikationslauf \code{cli109} deckt allerdings
nur Teil~1 ab: er lief, bevor Ihre zweite Runde bei uns bearbeitet war. Die Punkte
aus Teil~2 sind einzeln an echten Daten gepr\"uft, aber noch nicht in einem
vollst\"andigen Lauf. Der folgt als \code{cli110}.

\code{cli109}, 528 QC-Reports:

\begin{center}
\begin{tabular}{lrr}
\toprule
& \textbf{cli108} & \textbf{cli109} \\
\midrule
HIGH   & 1   & 1 \\
MEDIUM & 758 & 220 \\
TIME001 / TIME003 & \multicolumn{2}{c}{0 in cli109} \\
\bottomrule
\end{tabular}
\end{center}

Das verbleibende HIGH ist \code{basin} und stammt aus einer Regression in
\code{cc-plugin-wcrp} (\#66), nicht aus unseren Daten.

Neben den Review-Punkten haben wir eine eigene Suche nach Quelle-gegen-DReq-Fehlern
gemacht. Anlass war \code{mrsol}: dort zeigte sich, dass ein einzelner Ausgabestrom
mehrere DReq-Eintr\"age gleichzeitig falsch bedienen kann und \emph{kein} QC-Check das
findet. Diese Suche hat sieben weitere Punkte gefunden, darunter zwei mit falschen
Werten und einen, bei dem eine Datei gar nicht mehr lesbar war.

Umfang seit dem letzten Commit vor dem Review (\\code{b42279a2}): 16 Commits,
22 Dateien, +2291 / -362 Zeilen.

\section{Teil 1 Ihres Reviews: Antwort auf die Punkte}

\subsection*{1hr, \code{hfls}: \code{lon\_bnds}/\code{lat\_bnds} fehlen, keine \code{time}-Variable}
\ok\ Das war kein Metadatenfehler. Der Shard lief in die SLURM-Walltime, und ein
abgebrochener Schreibvorgang hinterl\"asst eine Header-only-Datei statt eines Fehlers.
Der Tier \code{extra\_atm} hat jetzt 6\,h. \code{cli109}: 8760 Zeitschritte,
\code{lat\_bnds} und \code{lon\_bnds} vorhanden.

Der Fehlermodus ist unangenehm, weil er lautlos ist. Wir haben deshalb auch
\code{cap7\_atm} vorsorglich auf 6\,h gesetzt, nachdem dort der Input um Faktor 26
gewachsen ist (siehe \code{hurs} im Abschnitt zur eigenen Nachsuche).

\subsection*{1hr, \code{hfls}: \code{\_QuantizeBitGroomNumberOfSignificantDigits}}
\ok\ Danke f\"ur den Hinweis auf CF 1.12. Umgesetzt wie in der Konvention:
ein Container \code{quantization\_info} mit \code{algorithm = "bitgroom"} und
\code{implementation = "netCDF-C version 4.9.3-development"}, dazu pro Variable
\code{quantization = "quantization\_info"} und \code{quantization\_nsd = 5}.
Das f\"uhrende Unterstrich-Attribut wird jetzt entfernt, das verstie\ss\ ohnehin
gegen CF \S2.3.

\subsection*{1hr, \code{rsus}: Repacking-Checks}
\offen, bewusst. Wir folgen Ihrem Vorschlag und lassen das vor der Publikation
durch \code{cmip7-repack} laufen.

\subsection*{3hr, \code{mrsol}: erster Zeitschritt (01:30) fehlt}
\ok, aber die Ursache lag tiefer, und der Hinweis war deshalb sehr n\"utzlich.

Ein einziger 3-st\"undiger XIOS-Stream (\code{operation="instant"}, Feld: oberster
Meter) bediente drei DReq-Eintr\"age, die sich sowohl in der Tiefe als auch in der
temporal shape unterscheiden. Zwei Folgen:

\begin{itemize}
\item \code{tavg-d100cm} bekam instantane Samples und konnte damit kein Mittel sein.
      Das war der von Ihnen gesehene fehlende erste Zeitschritt.
\item Beide \code{d10cm}-Eintr\"age lasen still das 100-cm-Feld, publizierten also
      1\,m Bodenfeuchte unter einem 10-cm-Label.
\end{itemize}

Der Tiefenfehler ist an cli108 belegbar: 3hr \code{d10cm} hatte ein Mittel von
72{,}6 kg\,m$^{-2}$ gegen 6{,}9 bei der echten 10-cm-Monatsregel, Verh\"altnis
10{,}5 und damit genau das Tiefenverh\"altnis 1{,}00\,m / 0{,}10\,m. Das sind falsche
Werte, kein Metadatenproblem, und kein QC-Check schl\"agt an.

Die drei Regeln sind vorerst deaktiviert. In \code{esm\_tools} sind drei getrennte
Streams erg\"anzt (100\,cm gemittelt, 10\,cm instantan, 10\,cm t\"aglich gemittelt),
sodass k\"unftige L\"aufe inklusive LR sie reaktivieren k\"onnen.

\subsection*{Allgemein: fehlende scalar coordinates}
\ok\ Sie hatten recht, die fehlten durchg\"angig. Wir lesen jetzt
\code{CMIP7\_coordinate.json} und schreiben beide Sorten:

\begin{itemize}
\item numerisch mit \code{value}: \code{height2m}, \code{height10m}, \code{sdepth10cm}, \dots
\item character: \code{typetree}, \code{typec3crop} und weitere Tile-/Area-Typen,
      \code{out\_name} meist \code{type}
\end{itemize}

Beide werden in \code{coordinates} der Datenvariable referenziert.
\code{cli109}, \code{sfcWind}: \code{coordinates = "time lat lon height"}.

\subsection*{Allgemein: globales Attribut \code{experiment}}
\ok\ Enth\"alt jetzt die CV-Beschreibung, \"uber \code{esgvoc} aufgel\"ost, statt
noch einmal die \code{experiment\_id}.

\subsection*{Allgemein: CF-Checker, \code{coordinates} als globales Attribut}
\ok\ Entfernt. Eine Warnung f\"ur andere Gruppen mit demselben Problem: der
naheliegende Weg \code{ds.encoding["coordinates"] = None} auf Dataset-Ebene wirkt
\emph{nicht}. Wir haben das empirisch gepr\"uft, xarray schreibt das Attribut
trotzdem. Nur \code{reset\_coords} auf den Bounds-Koordinaten hilft. Auf
\emph{Variablen}-Ebene funktioniert \code{encoding["coordinates"] = None} dagegen
sehr wohl, was die Sache verwirrend macht.

\subsection*{Allgemein: TIME001 squareness bei tavg}
\ok\ In \code{cli109} null Treffer, ebenso TIME003. Der Rest kam aus den
inzwischen gemergten Upstream-PRs plus unserer Korrektur des Dateinamen-Tokens auf
\code{time[0]}/\code{time[-1]} nach CMIP7 Appendix 1.

\section{Teil 2 Ihres Reviews: Antwort auf die Punkte}

Alle vier Punkte sind umgesetzt und einzeln an echten Daten gepr\"uft. Sie sind
noch nicht in \code{cli109} enthalten, das lief vorher.

\subsection*{\code{string} statt \code{character} bei Koordinatenvariablen}
\ok\ Und der Hinweis war wichtiger, als er aussah. Wir schreiben jetzt
\code{char sector(<out\_name>, strlen)}, gepr\"uft gegen publizierte
MPI-ESM1-2-LR-Dateien. \code{cSoilPools} war auf allen drei Punkten falsch und ist
jetzt \code{char sector(type, strlen)}, \code{long\_name = "Soil Pools"}, ohne
\code{units}. Bei \code{landCoverFrac} hei\ss t die Dimension jetzt \code{type} statt
\code{vegtype}; dass der \code{out\_name} zur Dimension wird, hatten wir zuerst
anders entschieden, Ihre Beschreibung war eindeutig.

Der Grund, warum das mehr als eine Konvention ist: \code{cli109} enth\"alt f\"unf
Dateien, die netCDF-C \emph{gar nicht} \"offnen kann.

\begin{quote}
\code{cropFrac}, \code{shrubFrac}, \code{baresoilFrac}, \code{treeFrac},
\code{grassFrac} (yr) \\
\code{RuntimeError: NetCDF: HDF error}
\end{quote}

Ihre skalare \code{type}-Koordinate ging als \code{NC\_STRING} raus, und h5py meldet
eine korrupte creation property list auf dem vlen-Dataset. Der Compliance-Checker
meldet dieselben Regeln mit \code{high=0 medium=0 low=0}. Wir hatten das zun\"achst
isoliert nachgestellt und dabei \emph{keinen} Fehler gesehen; er tritt erst im
vollst\"andigen Schreibpfad auf. Als Hinweis f\"ur die QC: eine Datei, welche die
Bibliothek nicht \"offnen kann, sollte kein sauberes Ergebnis bekommen.

\subsection*{\code{axis = X/Y} auf unstrukturierten Gittern}
\ok\ Entfernt, und Sie haben die Begr\"undung schon geliefert: \code{nod2} ist eine
Indexdimension. Wir setzen \code{axis} jetzt nur noch auf echten
Koordinatenvariablen, deren einzige Dimension ihr eigener Name ist.
\code{CMIP7\_grids.json} f\"uhrt bei \code{latitude}/\code{longitude} ebenfalls kein
\code{axis}. Betroffen waren 437 von 540 Dateien. Das \code{coordinates}-Attribut auf
\code{lat\_bnds}/\code{lon\_bnds} (274 Dateien) ist ebenfalls weg.

\subsection*{\code{olevel} ist nur ein Platzhalter}
\ok\ Aufgel\"ost zu \code{depth\_coord}, \code{out\_name} \code{lev}. Damit kamen
gleich drei Dinge in Ordnung: der Name, die von \code{must\_have\_bounds: yes}
verlangten \code{lev\_bnds}, und die Attribute. Letztere kamen bisher aus dem Modell
(\code{units = "meter"}, ein Nicht-CF-Attribut \code{name = "nz"}); jetzt gewinnt das
CV. Die fehlenden Bounds hatten eine eigene Ursache: der Bounds-Schritt l\"auft vor
dem Umbenennen und kannte den FESOM-Namen \code{nz1} nicht.

\subsection*{\code{alevel}: parametrische Koordinate mit \code{formula\_terms}}
\ok\ Gew\"ahlt ist \code{alternate\_hybrid\_sigma}, \code{p = ap + b*ps}. Die Dateien
enthalten jetzt \code{lev}, \code{lev\_bnds}, \code{ap}, \code{b}, \code{ap\_bnds},
\code{b\_bnds}, \code{ps} und die \code{formula\_terms}.

Die A/B-Koeffizienten stehen nicht im Modell-Output; sie kommen aus der
ECMWF-Vertikaltabelle f\"ur L137, die zur Diskretisierung geh\"ort und mit den
Modell-Inputdaten ausgeliefert wird. \code{ps} muss Gitter \emph{und} Zeitachse der
Variablen teilen: das cmorisierte \code{ps} liegt auf dem regul\"aren Regrid und
scheidet damit aus, also lesen wir den nativen 6-st\"undlichen Strom und mitteln ihn
\"uber die Ausgabeintervalle.

Zwei Fehler dabei sind uns nur aufgefallen, weil wir Feld f\"ur Feld gegen eine
publizierte MPI-ESM-Datei verglichen haben, nicht weil QC etwas gesagt h\"atte:
\code{lev} bekam zwischenzeitlich die Ozean-Metadaten (\code{standard\_name = depth},
\code{units = m}) auf 137 Modelllevel, und \code{ap}/\code{b} wurden mit BitGroom auf
5 signifikante Stellen quantisiert, also verlustbehaftet genau auf den Koeffizienten,
aus denen die Vertikalachse rekonstruiert wird. Beides ist behoben, \code{ap} und
\code{b} sind bitgenau gegen die Tabelle.

\subsection*{Nebenbefund aus Ihrem Hinweis auf \code{CMIP7\_grids.json}: \code{grid\_label}}

Beim Nachsehen ist uns aufgefallen, dass 161 Regeln ein \code{grid\_label} trugen,
das ein anderes Gitter beschreibt, darunter 142 Dateien mit \code{g113}. Das ist ein
regul\"ares 0{,}25-Grad-Gitter mit 1036800 Zellen und geh\"ort nicht uns.

Ursache: \code{grid\_label} stand einmal pro Tier im \code{inherit}-Block, und ein
Tier, dessen Regeln nativen und regemappten Strom mischen, stempelte beiden dasselbe
Label auf. Da \code{grid\_label} in Pfad \emph{und} Dateinamen steht, l\"asst sich das
nachtr\"aglich nicht reparieren, deshalb vor \code{cli110}.

Alle vier Gitter, die wir tats\"achlich schreiben, waren bereits registriert:

\begin{center}
\begin{tabular}{lll}
\toprule
\textbf{Gitter} & \textbf{Zellen} & \textbf{Label} \\
\midrule
OIFS TCo319 reduziert gau\ss sch & 421120 & \code{g122} \\
OIFS Regrid 640$\times$1296 & 829440 & \code{g129} \\
FESOM DARS Knoten & 3146761 & \code{g130} \\
FESOM DARS Elemente & 6229020 & \code{g132} \\
\bottomrule
\end{tabular}
\end{center}

Zwei weitere Zellmengen sind echte Teilmengen des TCo319-Gitters (LPJ-GUESS-Landpunkte
107132, Region 30S--90S 94572) und behalten \code{g122}. Beide hatten allerdings
\"uberhaupt keine Zellgrenzen, waren also gar nicht als TCo319-Zellen ausgedr\"uckt;
das ist jetzt korrigiert, 173 Dateien.

Offen ist nur, welches \code{grid\_label} eine Datei tragen soll, die gar kein
horizontales Gitter hat (globale Mittel, Becken-Schnitte, Profile; bei uns 36
Dateien). CMIP6 hatte \code{gm}, im numerischen Schema finden wir keine Entsprechung,
und \code{g999} hat eine leere Beschreibung. Falls Sie wissen, was hier vorgesehen
ist, w\"are das die letzte offene Frage aus diesem Block.

\section{Eigene Nachsuche (unabh\"angig von Ihrem Review)}

Ein Stream, der mehrere DReq-Eintr\"age bedient, kann diese gleichzeitig falsch
bedienen, und QC sieht davon nichts. Wir haben deshalb systematisch die Quelle jeder
Regel gegen ihren \code{compound\_name} gepr\"uft. Sieben Befunde, alle verifiziert:

\begin{enumerate}
\item \textbf{18 \code{*Lut}-Regeln ohne \code{landuse}-Achse.} Publiziert wurde nur
die Spalte \code{psl}. Die DReq-Dimensionen sind
\code{longitude latitude landuse time}. Alle vier Kacheln standen in den
Quelldateien, es war reine Loader-Arbeit. Jetzt 4 Werte in DReq-Reihenfolge.

\item \textbf{\code{landCoverFrac} ohne \code{vegtype}-Achse.} Die 44 PFT-Spalten
wurden zu einem Skalar aufsummiert. Jetzt 7 Werte. Die Zuordnung stammt aus den
Instruktionsdateien des Laufs selbst (\code{global.ins}, \code{arctic.ins},
\code{wetlandpfts.ins}, \code{crop\_n.ins}, \code{landcover.ins}), nicht aus den
PFT-K\"urzeln. Nicht zuordenbare Spalten l\"osen jetzt einen Fehler aus, statt still
wegzufallen. \code{CLM}, \code{pCLM}, \code{pmoss} und \code{Bare\_soil} haben keine
Entsprechung auf der \code{vegtype}-Achse und fehlen daher, was im
\code{comment} der Variable steht.

\item \textbf{\code{mrsolLut} las die ganze Bodens\"aule} unter \code{d10cm}-Label.
Die passende 0--10-cm-Ausgabe lag ungenutzt daneben. Umgestellt. Gleiche Fehlerklasse
wie bei \code{mrsol}.

\item \textbf{\code{sivol} t\"aglich: 12 Monatswerte auf 335 Tagesschritte},
davon 323 NaN. Die richtige Eingangsgr\"o\ss e \code{m\_ice} liegt nur monatlich vor,
\code{h\_ice} und \code{a\_ice} aber t\"aglich, und \code{m\_ice = h\_ice} $\cdot$
\code{a\_ice}. Wir rekonstruieren das Tagesfeld daraus.

Die N\"aherung ist gemessen, nicht gesch\"atzt: in einem Testlauf schreibt FESOM
\code{sivoln} jeden Modellzeitschritt (87\,600 Samples), das ist FESOMs eigene
Eisvolumen-Diagnostik. Gegen die Rekonstruktion:

\begin{center}
\begin{tabular}{lrrr}
\toprule
$10^3$\,km$^3$ & min & max & Mittel \\
\midrule
FESOM \code{sivoln}, auf Tage gemittelt & 10{,}911 & 29{,}497 & 20{,}585 \\
Rekonstruktion \code{h\_ice} $\cdot$ \code{a\_ice} & 10{,}921 & 29{,}524 & 20{,}605 \\
\bottomrule
\end{tabular}
\end{center}

Abweichung im Mittel $+0{,}094\,\%$, maximal $0{,}107\,\%$, Korrelation
$1{,}000000$. Der Rest ist die sub-t\"agliche Kovarianz zwischen Dicke und
Konzentration, die das Produkt zweier Tagesmittel verliert. \code{cli109}:
365 Zeitschritte, keine NaN. Ein t\"aglicher \code{m\_ice}-Stream ist in
\code{esm\_tools} erg\"anzt und macht die N\"aherung f\"ur sp\"atere L\"aufe
\"uberfl\"ussig.

\item \textbf{\code{sistressave} und \code{sistressmax}: Monatsmittel als \code{tpt}
gelabelt.} Beide DReq-Eintr\"age wollen einen Momentanwert zur Monatsmitte, FESOM
schreibt \code{sgm11}/\code{sgm12}/\code{sgm22} als Monatsmittel. Bei
\code{sistressmax} ist es mehr als ein Labelproblem: die Invariante
$\sqrt{((\sigma_{11}-\sigma_{22})/2)^2 + \sigma_{12}^2}$ ist nichtlinear in den
Komponenten, aus Monatsmitteln berechnet ist sie nicht einmal das Monatsmittel der
Invariante. Ausgabe pro Zeitschritt w\"are der einzige Weg und ist bei HR mit
grob 1{,}5\,TB pro Jahr und Komponente nicht machbar. Beide Regeln deaktiviert,
instantane Streams in \code{esm\_tools} erg\"anzt.

\item \textbf{\code{hurs} Tagesmaximum und -minimum aus Tagesmitteln.} Es fehlte
schlicht eine max/min-Reduktion. Die gibt es jetzt, gespeist aus Stundenwerten.
Das ist eine N\"aherung gegen\"uber Extrema auf Modellzeitschritt-Ebene, aber
konsistent mit dem, was viele CMIP7-Modelle liefern werden.

\item \textbf{\code{emibb*}: \code{fFireAll} oder \code{fFireNat}?} Die DReq gibt
keine Beschreibung. Dom\"anenentscheidung, best\"atigt als \code{fFireAll}
inklusive Landnutzungsfeuer. Das bisherige Verhalten war richtig, dokumentiert,
damit die Frage nicht wieder aufkommt.
\end{enumerate}

\subsection*{Nebenbefund: QC meldet eine unlesbare Datei als sauber}

Beim Bau der \code{landuse}-Achse haben wir die Labels zun\"achst als
Python-Strings abgelegt. xarray schreibt daraus eine \code{NC\_STRING}-Variable, und
die resultierende Datei lie\ss\ sich von netCDF-C \"uberhaupt nicht mehr \"offnen:
\code{ncdump} und \code{netCDF4.Dataset} scheiterten beide mit
\code{NetCDF: HDF error}, h5py meldete eine korrupte creation property list
(\code{bad global heap collection signature}) auf dem vlen-Dataset.

Der Compliance-Checker meldete dieselbe Datei mit \code{high=0 medium=0 low=0}.

Wir schreiben die Achsen jetzt als \code{char sector(landuse, strlen)}, also so, wie
CMOR es tut, gepr\"uft gegen publizierte MPI-ESM1-2-LR-Dateien. Das behebt es
vollst\"andig. Der Hinweis k\"onnte f\"ur die QC allgemein interessant sein: eine
Datei, die die Bibliothek nicht \"offnen kann, sollte kein sauberes Ergebnis
bekommen.

\section{Verbleibende Befunde in cli109}

\begin{longtable}{rll}
\toprule
\textbf{Anzahl} & \textbf{Check} & \textbf{Einordnung} \\
\midrule
\endhead
1   & \code{basin} HIGH   & \upstream, \code{cc-plugin-wcrp} \#66 \\
162 & \code{ATTR004}      & Registry-Erwartung f\"ur \code{long\_name} / \code{cell\_methods} \\
36  & \code{ATTR001}      & \code{cell\_measures} auf \code{areacella} und auf hemisph\"aren- \\
    &                     & integrierten Skalaren, die keins haben \\
17  & CF \S2.4            & Dimensionsreihenfolge \\
5   & CF \S6.1            & Labels, \code{basin}, inzwischen behoben (s.\,u.) \\
\bottomrule
\end{longtable}

Die f\"unf \S6.1-Befunde haben wir nach dem Lauf noch aufgekl\"art. Der Checker
meldete auf \code{hfbasin} und \code{sltbasin}:

\begin{quote}
\code{6.1.1 'atlantic\_arctic\_oceanindian\_pacific\_oceanglobal\_ocean'
specified by 'basin' is not a valid region}
\end{quote}

Die drei Labels laufen zu einem String zusammen. Das ist dieselbe Ursache wie
beim \code{landuse}-Nebenbefund oben: die Koordinate enthielt Python-Strings,
xarray schrieb eine \code{NC\_STRING}-Variable, und der Checker liest die Achse
dann als einen einzigen Wert statt als drei. Jedes Label f\"ur sich ist im CV
g\"ultig. Wir schreiben die Achse jetzt als \code{char sector(basin, strlen)},
mit \code{standard\_name = region} und \code{long\_name = "Ocean Basin"} aus dem
CMIP7-CV, gepr\"uft gegen publizierte MPI-ESM1-2-LR-Dateien (\code{Omon.hfbasin}).
Die Verifikation \"uber die drei betroffenen Regeln kommt mit der n\"achsten Runde.

Bei \code{ATTR004} ist mindestens ein Teil upstream: f\"ur
\code{seaIce.sivol.tavg-u-hm-u.day.sh} erwartet die Registry
\code{"Sea-Ice Volume North"}, w\"ahrend die DReq-Tabelle
\code{"Sea-Ice Volume South"} vorgibt. Denselben Befund gab es schon in cli108 bei
\code{sivol\_south} und \code{sivol\_south\_day}. Wir gehen die Liste noch durch und
melden, was upstream geh\"ort.

\code{rsus} Repacking l\"auft wie besprochen vor der Publikation \"uber
\code{cmip7-repack}.

\section{Offene Frage an Sie}

Nur \code{mrsol} wurde urspr\"unglich von Hand gepr\"uft, weil Ihr Review darauf
zeigte. Unsere Nachsuche hat eine Stichprobe abgedeckt, nicht alles. Andere
Variablen, die sich einen Stream \"uber mehrere DReq-Eintr\"age teilen, k\"onnten
dieselbe Fehlerklasse haben, und sie ist f\"ur QC unsichtbar. Wir planen vor der
Publikation einen Durchlauf, der f\"ur jede Regel \code{compound\_name} gegen
\code{operation} und Felddefinition des gelesenen Streams pr\"uft. Falls es so etwas
schon gibt oder bei DKRZ geplant ist, w\"are das gut zu wissen.

\end{document}
